% Interspeech version (4 pages + references) of ../main.tex, the full paper. Build with XeTeX: see ../README.md.
\documentclass{Interspeech}
\usepackage{fontspec}
\usepackage{xeCJK}
\setmainfont{texgyretermes}[Extension=.otf,UprightFont=*-regular,BoldFont=*-bold,ItalicFont=*-italic,BoldItalicFont=*-bolditalic]
\newfontfamily\ipafont{Charis SIL}[Scale=0.92] % IPA, including the pitch-accent letters U+A71B / U+A71C
\setCJKmainfont{Hiragino Mincho ProN}
\usepackage{booktabs}

% YomiScript and IPA examples: copied exactly from the spec and the result files
\newcommand{\ys}[1]{\mbox{\ipafont #1}}

\title{The LLM as the Front End: Context-Aware Pronunciation for Small Browser Text-to-Speech}
\author[affiliation={1}]{Anonymous}{Author}
\address{$^1$Anonymous Institution}
\email{anonymous@example.com}
\keywords{text-to-speech, grapheme-to-phoneme, homographs, large language models, Japanese, pronunciation notation}

\begin{document}
\maketitle

\begin{abstract}
Japanese text-to-speech must often choose a reading from context: homographs, counters, dates and acronyms. We let
an LLM, given the text and its context, decide the reading, and pass it to the synthesizer in YomiScript, a notation
built on IPA and BCP~47, so that an adapter for Kokoro-82M needs no text analysis and runs in a browser. On 130
sentences built so that only the context decides (about 50\% expected without it), three LLMs read the
target word right in 48.5\% without context and 97.7\% with the preceding paragraph. On an independent 3{,}000-sentence
benchmark with frozen prompts, two LLMs' readings had a kana error rate of 0.25--0.27\% written as kana, against
1.05\% for pyopenjtalk; written directly as YomiScript they had 0.63--0.87\%, or 1.50--3.47\% counting the utterances
the synthesizer rejects. The LLM should write Japanese readings as kana, converted by table; the notation carries what
kana cannot, accent and other languages, measured by pitch and ASR, not yet by listening.
\end{abstract}

\section{Introduction}

Small open TTS models such as Kokoro-82M~\cite{kokoro82m} can run in the browser tab that shows an LLM chat. What does
not move easily is the front end that decides how each word is read. Kokoro's Japanese front end
(misaki~\cite{misaki}, MeCab~\cite{kudo2004mecab} with UniDic~3.1.0~\cite{den2008unidic}) needs a 481\,MB connection table
and a 243\,MB dictionary. Smaller front ends fit (OpenJTalk~\cite{openjtalk} compiled to WebAssembly is about
19\,MB~\cite{piperplus}), but they read the sentence, not what it means in its context: 一日 is \emph{tsuitachi} (the
first of the month) or \emph{ichinichi} (one day), 上手 is \emph{jōzu} (skilful) or \emph{kamite} (stage left, the audience's right), and
which one is meant is often decided outside the sentence.

An LLM can use the meaning of the text and any context supplied with it. We let an LLM decide the reading and hand
the TTS a resolved pronunciation through an interface, so the synthesizer needs no reading analysis. What crosses the
interface is YomiScript; how the Yomite produces it is its own choice, and our results favour an LLM that writes kana,
a table that turns kana into IPA, and an optional dictionary step that adds accent from the original text. Our
contributions: (1)~\textbf{YomiScript}, a TTS notation built on IPA~\cite{ipa1999handbook} and BCP~47~\cite{rfc5646}
that separates the \emph{Yomite} (text to pronunciation: an LLM, rules, a person) from the \emph{Katarite}
(pronunciation to sound), and a Kokoro Katarite that runs in the browser and, given the original front end's
pronunciation, reproduces the original Kokoro's model input exactly;
(2)~\textbf{context experiments}: 65 homographs in sentences we wrote to be ambiguous alone, read with and without
a disambiguating paragraph or conversation, and by the LLM that wrote them, with the main comparison stated in advance;
they test whether LLMs use such context, not how often natural text needs it;
(3)~an evaluation on an \textbf{independent benchmark}~\cite{koriyama2026g2p} with frozen prompts, and a
\textbf{container comparison}: the same LLMs writing YomiScript, kana, SSML and sparse kana, and what kana cannot carry
(accent and English), measured in the synthesized pitch and by ASR, not by listening.

\section{Related work}

OpenJTalk and pyopenjtalk~\cite{pyopenjtalk} read Japanese with MeCab and NAIST-JDIC, and accent can be predicted by
multi-task learning~\cite{park2022marine}. Contextual embeddings disambiguate English homographs~\cite{nicolis2021homograph},
and pre-trained language model features improved Japanese polyphone disambiguation and accent prediction by about six
points (5.7 for polyphones, 6.0 for accent)~\cite{hida2022polyphone}; such a model is the baseline we have not run. LLMs have been put into the front end.
ICKR has an LLM choose the sense of English homographs and a dictionary give the phonemes~\cite{han2024ickr}, the same
division of labour our results favour for Japanese; PolyNorm normalises text with an LLM in eight
languages~\cite{wong2025polynorm}; and Koriyama benchmarks over 30 LLMs on Japanese G2P (3{,}000 sentences), the best
at a kana error rate of 0.52\% with an LLM analysis converted to kana by rule, against 1.03\% for
OpenJTalk~\cite{koriyama2026g2p}. These work on the sentence. We give the LLM the context the text comes with and
measure how much it matters, hand over a notation not tied to one engine's input, and keep no reading analysis in the
Kokoro Katarite. SSML~\cite{ssml11} has \texttt{phoneme} and \texttt{xml:lang} and can carry the same IPA; YomiScript
differs as a compact, pronunciation-first convention in which every word is resolved, accent and phrasing have their
IPA symbols, and the original text is an optional layer.

\section{YomiScript and the Katarites}

An utterance is a BCP~47 tag and IPA between slashes, ending in a fall \ys{↘} or rise \ys{↗}; a \ys{\{en-US …\}}
section switches language inside it. Japanese is written mora by mora with the specification's mora table (ザ as
\ys{dza}, ニ as \ys{ɲi}); \ys{ꜛ} marks where the pitch rises and \ys{ꜜ} falls after the accent nucleus, so in Tokyo
Japanese 雨が is \ys{aꜜmeɡa} and 飴が \ys{aꜛmeɡa}; \ys{|} and \ys{‖} mark phrase and breath-group
boundaries, and optional pitch values may follow a mora. The parser rejects symbols outside the inventory (a Latin
\ys{r}, an unpaired mark). A text layer \ys{<t:…>} closes each span with the text it reads; it never
changes the segmental reading (the parser removes it), lets a client show the text as it is spoken, and lets a later
step add the dictionary accent of the text where its reading agrees (OpenJTalk compiled to WebAssembly with a 28\,MB
dictionary: a front end again, but one that only adds pitch). An LLM output for the hard-set sentence 一日から三日まで休みです
(``closed from the 1st to the 3rd''):

\begin{center}\footnotesize
\ys{ja /tsɯitatɕi<t:一日> kaɾa<t:から> mikka<t:三日>}\\
\ys{made<t:まで> jasɯmi<t:休み> desɯ<t:です。> ↘/}
\end{center}

The \emph{Kokoro Katarite} maps YomiScript IPA to Kokoro's phonemes and runs its three ONNX
parts~\cite{onnxruntime} in the browser. Given YomiScript written from Kokoro's own front end, it hands the model the
same phoneme string, ids, voice entry, durations and F0 as the original pipeline on 23/23 English and 50/50 Japanese
chunks (Table~\ref{tab:parity}), in a mode that renders Japanese pauses as the original front end does. With fp32
weights the audio's correlation with the original is that of a second run of the original, whose decoder draws
random numbers. The browser's fp16 weights are close on average (0.985, 0.983) but their worst chunks fall to 0.85
and 0.80, against 0.99 and 0.98 for a second run; we have not listened for the difference.

We measured the browser path on one laptop (Apple M2 Max) in headless Chromium, with the models served over
loopback: the fp16 model (163\,MB; 186\,MB transferred on a cold load) was ready 1.7\,s after opening the page and
spoke its first 2\,s sentence at 2.8\,s; the real-time factor was 0.54 on four WebAssembly threads and 0.06 with the
decoder on WebGPU (with the encoder on WebGPU too, the predicted durations came out wrong); the renderer process held about 1.7\,GB when ready. A second Katarite for Supertonic~3~\cite{supertonic3}, which reads characters, gives kana from the
YomiScript only for written forms with several readings.

\begin{table}[t]
\centering\footnotesize
\caption{The Kokoro Katarite against the original Kokoro pipeline, given the original front end's pronunciation as
YomiScript: identical model input, and the waveform's correlation with the original's (mean / min).}
\label{tab:parity}
\setlength{\tabcolsep}{3pt}
\begin{tabular}{lcc}
\toprule
& English (23) & Japanese (50)\\
\midrule
Input, durations, F0 & identical & identical\\
Original, second run & 0.9943 / 0.9924 & 0.9918 / 0.9762\\
Katarite, fp32 & 0.9944 / 0.9888 & 0.9914 / 0.9773\\
Katarite, fp16 (browser) & 0.9854 / 0.8545 & 0.9833 / 0.7980\\
\bottomrule
\end{tabular}
\end{table}

\section{Experimental setup}

\textbf{LLMs.} Claude Code (\texttt{claude-opus-5-5}), Codex (\texttt{gpt-6.1-sol}) and Antigravity with
Gemini (\texttt{gemini-4-argon-medium}) through their CLIs (29 October to 1 November 2026) with no tools and the CLIs' default reasoning (they expose no
sampling controls), an answer without JSON asked once more, given the YomiScript instructions our tools ship (the
prompts are released with the code).
Readings are taken from the YomiScript (as the phonemes Kokoro is given, or through the text layer for a target word)
and compared with katakana gold, merging spellings that do not change the reading; for a target word only its
reading is compared, whether or not the rest of the utterance parses. The project-created sets and their readings
were reviewed by one native speaker of the project.

\textbf{Context.} For 65 homographs we wrote 130 items: both items of a word share one short target sentence that is
ambiguous alone (「それは生花です。」), and each has a preceding paragraph (like a web page or a document) and a user's
message that fix one reading without containing the word or its reading (checked by script). Without information about the intended sense, about half right is the
expected best. Calls held 10 items, never two of the same word. The LLMs wrote YomiScript for the target sentence as external text with and
without the paragraph, and as a fixed chat answer with and without the user's message; as authors they also answered
each message in at most 15 characters, and every LLM was to read every author's answer with and without the message;
Gemini's quota ran out during these reader calls, so the author comparison uses the 210 of 390 answers every reader
returned. The
primary comparison (external text, with minus without the paragraph, pooled) was stated before the run.

\textbf{Containers.} On the 70-sentence hard set (35 homograph, 12 acronym, 12 number and 11 date sentences), the
same three LLMs wrote YomiScript, kana (as in~\cite{koriyama2026g2p}), SSML with an IPA \texttt{phoneme} on every word,
and sparse kana (readings after ambiguous words only, the rest read by pyopenjtalk), three runs each, one call per set.
The prompt was revised twice on this set, for notation only, so these are development results.

\textbf{Independent benchmark.} JVS-nonpara-kana~\cite{koriyama2026g2p}: 3{,}000 sentences of the JVS
corpus~\cite{takamichi2020jvs} with hand-annotated kana, scored by kana character error rate (CER) with Koriyama's
normalisation (reproduced exactly on all 3{,}000 gold readings), merging ヲ/オ, ヂ/ジ and ヅ/ズ, which his gold
writes one way. Where a YomiScript answer did not parse, it was parsed word by word and the rejected words counted as
deleted, so this CER is more lenient than the Katarite, which rejects the utterance. The YomiScript and kana prompts were
used unchanged from the container experiment, 50 sentences per call, one run; YomiScript is turned into kana by
inverting the mora table of the specification; kana written by an LLM goes through that table the other way to
become YomiScript, so the kana row is also the CER of a Yomite that writes kana and converts it deterministically.
Gemini ran out of quota after the first 1{,}250 sentences of our seeded order. JVS texts are public, so they may be in the LLMs' training data.

\textbf{What kana cannot carry.} For these the prompt was changed to ask for accent marks and for \ys{\{en-US …\}}
sections. 44 words in 21 Tokyo minimal sets (雨/飴, 橋/箸/端, \dots) in sentences that fix
the word: the LLMs wrote accent marks (three runs), scored against OpenJTalk's accent; Kokoro (one voice) synthesized
the first run's sentences with the marks, with its own pitch, and from kana with a dictionary's guess, and we measured F0 per mora
(Praat~\cite{boersma2001praat}) at the end of each mora, where the rule pitch of the marks has settled. 16 Japanese
sentences with English meant as English (and 8 loanword controls): the
English stretches were transcribed with Whisper~\cite{radford2023whisper}: the YomiScript ones as synthesized
sections, the kana ones cut from the running speech, which favours YomiScript somewhat.

\section{Results}

\begin{table}[t]
\centering\footnotesize
\caption{Homograph read right in 130 items (65 words; ambiguous target sentences, ceiling about 50\% without
context), one run, 95\% interval from resampling words. Author rows: 210 answers every reader returned.}
\label{tab:context}
\setlength{\tabcolsep}{3pt}
\begin{tabular}{lc}
\toprule
Condition (3 LLMs pooled) & read right\\
\midrule
External text, sentence alone & 48.5\% [45.6, 51.0]\\
External text, with the preceding paragraph & \textbf{97.7\%} [95.4, 99.5]\\
Chat answer alone & 48.5\% [46.4, 50.5]\\
Chat answer, with the user's message & \textbf{98.2\%} [96.4, 99.7]\\
pyopenjtalk, target sentence & 49.2\% [47.7, 50.0]\\
\midrule
Author of a terse answer & 99.0\% [97.6, 100]\\
Same LLM, with the user's message & 99.0\% [97.6, 100]\\
Other LLMs, with the user's message & 98.6\% [97.4, 99.5]\\
Same LLM, answer alone & 75.2\% [67.6, 82.4]\\
\bottomrule
\end{tabular}
\end{table}

\textbf{Context decides the reading} (Table~\ref{tab:context}). With the preceding paragraph the LLMs read 49.2 points
more items right than from the sentence alone [+46.2, +52.3]; 64 words were read better and none worse (sign test
$p \approx 10^{-19}$), with similar point estimates for each LLM (+46.9 to +51.5; the lowest interval bound is +42.3). A user's message worked as well as a
paragraph. The text need not come from the LLM, but the LLM that writes it has the context in the same call: on the
answers every reader returned, the author and a reader given the message had equal observed accuracy, no difference detected (0.0 points
[$-$1.9, +1.9]), and a reader of the answer alone 23.8 points worse (32 words to 1). Only 5 of 390 terse answers
repeated a word of the message. Context failed on a few words (河岸 meant as \emph{kawagishi} was read \emph{kashi}
almost everywhere). Two readings listed in dictionaries were accepted after seeing the misses; without them the
primary difference is +48.2 [+44.9, +51.5]. An earlier, smaller run (25 words, answers not limited in length) gave
the same picture: 98.0\% for the author, 97--98\% for a reader given the conversation and 92.0\% given the answer only.

\begin{table}[t]
\centering\footnotesize
\caption{The hard set (70 sentences) written in four containers by three LLMs, three runs each, pooled: sentences and
target words (of 69) read right, malformed or missing items per call, characters per sentence (input 12.3) and seconds per call.}
\label{tab:containers}
\setlength{\tabcolsep}{2.5pt}
\begin{tabular}{lccccc}
\toprule
Form & sent.\,\% & words\,\% & malf. & chars & s\\
\midrule
YomiScript & 98.1 & 99.4 & 0.1 & 75 & 68\\
Kana$^{a}$ & 88.9 (100) & 88.9 (100) & 7.8 & 16 & 26\\
SSML, every word & 96.8 & 98.1 & 1.6 & 307 & 107\\
Sparse kana & 98.4 & 99.8 & 0.1 & 20 & 24\\
pyopenjtalk & 78.6 & 81.2 & -- & -- & --\\
Kokoro's front end & 32.9 & 31.9 & -- & -- & --\\
\bottomrule
\end{tabular}\\[2pt]
{\scriptsize $^{a}$One Gemini run returned no JSON and counts as wrong; in parentheses, the eight runs that answered. Malformed: items per 70-item call.
The last two rows are deterministic, one run.}
\end{table}

\textbf{On the hard set} (Table~\ref{tab:containers}) the LLMs read 98.1\% of sentences right, against 78.6\% for
pyopenjtalk and 32.9\% for Kokoro's front end, which leaves acronyms in Latin letters and misreads dates (24日 as \emph{nijūyon-ka}). pyopenjtalk reads the acronyms and dates right; its gap is mostly homographs (21 of the 35 homograph sentences
right; 11 of 12 number sentences). Per LLM, YomiScript scored 96.2\% (Claude), 99.5\% (Codex) and 98.6\%
(Gemini); its misses were mostly slips in writing IPA (Latin \ys{u} for \ys{ɯ}, a lost \ys{ʲ}), not other readings.
\textbf{For Japanese readings alone, kana suffices}: kana was read right in every sentence of every
run that answered, at a fifth of the characters; sparse kana needs a front end for the rest (pyopenjtalk read 日本 as
\emph{nippon} in every run). SSML with one element per word split geminates across elements (\ys{moɾat} + た); with
one element per bunsetsu (and calls of 10, below, so the two changes are confounded) that failure disappeared and
its accuracy (96.2\%) was 1.2 points below YomiScript's in the same calls, at 2.6 times YomiScript's
characters. Because all 70 sentences went in one call, the two sentences of a homograph could inform each other; rerun
in calls of 10 with no shared word (Claude and Codex, three runs), the scores moved by under a point (YomiScript
97.4\%, kana 99.3\%; differences $-$0.5 [$-$2.9, +1.7] and $-$0.7 [$-$1.9, 0.0]).

\begin{table}[t]
\centering\footnotesize
\caption{What kana cannot carry. Top: Tokyo minimal pairs (44 words), the share of compared mora points (n) where
the difference in Kokoro's pitch between the members of a pair has the sign their accents predict; not a listening
result. Bottom: English
stretches in 16 Japanese sentences, words recognised by Whisper.}
\label{tab:kana}
\setlength{\tabcolsep}{3pt}
\begin{tabular}{lc}
\toprule
Accent source for Kokoro & points right (n)\\
\midrule
LLM marks (Claude, Codex, Gemini) & 92, 95, 95\% (20--24)\\
Dictionary marks on the text layer & 95\% (22)\\
Kokoro's own pitch (no marks) & 42\% (24)\\
Kana, accent guessed by a dictionary & 19\% (26)\\
\midrule
English stretch written as & words heard\\
\midrule
YomiScript \ys{en-US} section & 96--99\%\\
Kana (katakana English) & 65--75\%\\
\bottomrule
\end{tabular}
\end{table}

\begin{table}[t]
\centering\footnotesize
\caption{Independent benchmark (JVS-nonpara-kana, frozen prompts, one run): kana CER in \%, with 95\% intervals
from resampling sentences; ヲ/オ, ヂ/ジ and ヅ/ズ merged (Koriyama's gold writes を as オ; unmerged, the kana rows are
2.11--2.14 and pyopenjtalk 2.88, while YomiScript, converted with オ, does not change). Gemini: the first 1{,}250
sentences. Koriyama reports 0.52 for the best LLM and 1.03 for
OpenJTalk~\cite{koriyama2026g2p}.}
\label{tab:koriyama}
\setlength{\tabcolsep}{3pt}
\begin{tabular}{lcc}
\toprule
System & 3{,}000 sentences & first 1{,}250\\
\midrule
Kana, Claude & 0.27 [0.20, 0.34] & 0.35\\
Kana, Codex & 0.25 [0.20, 0.30] & 0.35\\
Kana, Gemini & -- & 0.47\\
YomiScript, Claude & 0.63 [0.53, 0.73] & 0.77\\
\quad strict, rejected = deleted & 1.50 [1.09, 1.96] & --\\
YomiScript, Codex & 0.87 [0.77, 0.99] & 0.88\\
\quad strict, rejected = deleted & 3.47 [2.77, 4.22] & --\\
YomiScript, Gemini & -- & 0.97\\
pyopenjtalk & 1.05 [0.92, 1.21] & 1.06\\
\bottomrule
\end{tabular}
\end{table}

\textbf{On independent text, the LLMs read better than the dictionary front end, and kana is the more accurate
container} (Table~\ref{tab:koriyama}). Our pyopenjtalk scores 1.05\%, consistent with Koriyama's 1.03\% for
OpenJTalk under his normalisation. Written as kana, Claude's and Codex's readings had about a quarter of pyopenjtalk's errors
(0.25--0.27\%; Koriyama's best LLMs, older models in another prompt setup, had 0.52\%); written as YomiScript, 60\% and
83\% of pyopenjtalk's, 0.36 and 0.62 points more than as kana [0.29, 0.44] and [0.53, 0.73]. Of the extra
error characters (315 for Claude, 546 for Codex), about three quarters are IPA the Katarite rejects (a Latin \ys{r},
a geminate split by a word space: 120 and 282), いう written as \emph{iu} where the gold has \emph{yū} (81 and 78
more than as kana; the kana prompt has an example for it, the YomiScript prompt none), morae with no kana (22, 25)
and 先生, which both prompts' examples spell \ys{seːseː} / セーセー (13 and 2 more than as kana); the rest
(79, 159) are other slips and readings. Part of this looks like prompt guidance. A post-hoc matched YomiScript prompt (the いう rule and a corrected 先生 example), run on all 3{,}000 sentences, nearly
removed those errors but scored 0.54\% and 0.82\%, closing 25\% and 8\% of the gap to kana, because more of its
utterances were rejected. Counting a rejected utterance as wholly deleted, as the Katarite would speak it, frozen
YomiScript scores 1.50\% and 3.47\% (1--3\% of utterances rejected), against 0.25--0.27\% for kana. With these prompts, a Yomite that writes
kana and lets a table turn it into YomiScript was the more accurate; the Katarite's input is the same.

\textbf{What kana cannot carry} (Table~\ref{tab:kana}). The LLMs placed the accent as OpenJTalk does in 96--98\% of
the 44 minimal-pair words over three runs (the one disagreement, 潮, is a word dictionaries disagree on). With the
marks, the pitch difference between the members of a pair had the predicted sign at 92--95\% of the compared points
(median 4.5--5 semitones), against 42\% with Kokoro's own pitch and 19\% from kana, whose accent a dictionary must
guess without the kanji (15 of 44 words right). The points are few (20--26 per condition, shared within minimal
sets); measured over the whole mora instead of its end, the expected rises and falls within a word
are right at only 71--81\%, as the rule pitch puts the peak of a first-mora accent late;
nobody has yet judged by ear whether the words are told apart. The dictionary on the text layer, which also supplies the gold, gives these marks by construction,
so for Japanese the accent need not come from the LLM; kana without the text loses it (sparse kana keeps the text, and
whole-sentence kana would need it aligned back). In the
code-switched sentences the LLMs put all 16 English stretches in \ys{en-US} sections (Gemini 15, one answer
rejected) and none of the 8 loanwords; Whisper recognised 96--99\% of those English words, against 65--75\% from
the same stretch said in kana.

\section{Discussion and limitations}

A dictionary front end chooses from the sentence; an LLM reads the meaning and uses as much context as it is given.
On self-contained sentences that was enough; on sentences ambiguous alone the surrounding text decided, whether it
came from elsewhere or from the conversation. The Yomite need not be the writer, but in an LLM application the writer
already has the context and can write the reading in the same call, at the cost of output tokens: YomiScript is
5.6--6.1 times the characters of Japanese text (2.7--3.2 without the text layer), and every result here comes from a
frontier model whose calls took tens of seconds. YomiScript is best seen as the interface between the parts: for
Japanese readings the LLM can write kana, converted by table, with the accent added from the text layer; the notation
earns its size where kana cannot go, in accent marks the synthesizer follows and in other languages.

\textbf{Limitations.} (1)~The sets are small and written by the project, possibly with LLM help; their readings were
reviewed by one native speaker of the project, not by independent annotators. (2)~The hard-set prompt was revised
twice on the same set (for notation) and the scoring merges were chosen while looking at outputs, so
Table~\ref{tab:containers} is a development result; the context experiment's main comparison was stated in advance,
but two readings were added to its scoring afterwards. (3)~One run for the context experiment, three for the
containers; Gemini's quota ran out during the context readers' calls, so the author comparison covers 210 of 390
answers. (4)~Only frontier LLMs were tested; smaller and local models were not. (5)~Accent and code-switching were
measured in pitch and by ASR with one voice, and nothing here was checked by listening. (6)~No context-aware Japanese G2P baseline~\cite{hida2022polyphone}; the independent benchmark was
run once with two LLMs (Gemini on 1{,}250 sentences), and the YomiScript prompt has notation gaps it exposed (いう).
(7)~The context items were built so that context decides; they show that LLMs use it, not how often natural text
needs it. A larger, independently annotated test with frozen
prompts, and listening tests, are the next step.

\section{Conclusion}

An LLM given the context can decide how text is read, and a notation built on IPA and BCP~47 carries that decision to
a small synthesizer with no front end of its own. In our experiments the context explained the gain, and no
effect of authorship was detected once a reader had the context. On independent text, LLM readings had about a quarter (as kana) to 60--83\% (as YomiScript) of a dictionary
front end's errors. With the prompts tested, Japanese readings were more accurate written as kana and turned into IPA by table; the notation
carries what kana cannot, accent and other languages, whose effect on listeners remains to be tested. These
conclusions hold for the frontier LLMs and the constructed items tested here.

\bibliographystyle{IEEEtran}
\bibliography{../references}

\end{document}
